% GENERATED FILE. Edit canonical lessons, then run npm run generate:books.
% canonical-source-hash: fnv1a64:07c51c549a2dcaf0
% canonical-lessons: MR-C457-khed, MR-C457-pascattap, MR-C457-hundka, MR-C457-halhalne, MR-C457-musmusne, MR-R457-review-450-457

\chapter{Regret and Tears}
\label{ch:mr-regret-and-tears}
\hlchaptermodality{457}

\begin{chapteropening}
\textbf{By the end of this chapter:} \emph{I can talk about regret, remorse and quiet tears.}

Complete the last lesson of chapter 457: I can talk about regret, remorse and quiet tears.
\end{chapteropening}

\section[\texorpdfstring{\mr{खेद} (khed)}{खेद (khed)}]{\mr{खेद} (khed) — regret}

\label{lesson:MR-C457-khed}



\begin{quote}
Before the new one: say the Marathi for to get flustered, then the Marathi for to be bashful.
\end{quote}

\subsection*{You'll want to know: \mr{खेद}}
\textbf{\mr{खेद}} — \emph{khed} — \textquotedblleft{}regret\textquotedblright{}.

\begin{quote}
\mr{त्याला} \mr{आपल्या} \mr{चुकीचा} \mr{खेद} \mr{वाटला}.
\end{quote}

\begin{quote}
\emph{tyālā āplyā cukīcā khed vāṭlā.}
\end{quote}

\begin{quote}
He regretted his mistake.
\end{quote}

Masculine.

Formal: \mr{खेद} \mr{वाटणे}, to regret, with the dative and the genitive of the thing.

\begin{cousinweb}
\textbf{\mr{खेद}} (\emph{khed}) is a tatsama, taken from Sanskrit \emph{kheda}, \textquotedblleft{}lassitude, pain\textquotedblright{} (CDIAL 3919); in Marathi it is regret.
\end{cousinweb}

\begin{grammarlens}[title={What you've built}]
One more word for describing what you lived and felt.
\end{grammarlens}

\subsection*{Your turn}
\begin{itemize}
\raggedright
  \item \emph{Say it:} \emph{khed}
  \item \emph{Say it:} \emph{khed}, once more
  \item \emph{Say it:} \emph{tyālā āplyā cukīcā khed vāṭlā.} — He regretted his mistake.
  \item \emph{Recall:} say the Marathi for to get flustered, then the Marathi for to be bashful, then say \emph{khed} again
\end{itemize}

\begin{tcolorbox}[breakable,colback=teal!4,colframe=teal!35!black,title={Before you move on}]
What does \mr{खेद} mean? (\textquotedblleft{}Regret\textquotedblright{}.) Say it once more.
\end{tcolorbox}
\section[\texorpdfstring{\mr{पश्चात्ताप} (paścāttāp)}{पश्चात्ताप (paścāttāp)}]{\mr{पश्चात्ताप} (paścāttāp) — remorse}

\label{lesson:MR-C457-pascattap}



\begin{quote}
Before the new one: say the Marathi for to be bashful, then the Marathi for regret.
\end{quote}

\subsection*{You'll want to know: \mr{पश्चात्ताप}}
\textbf{\mr{पश्चात्ताप}} — \emph{paścāttāp} — \textquotedblleft{}remorse\textquotedblright{}.

\begin{quote}
\mr{रागात} \mr{बोलल्याचा} \mr{त्याला} \mr{नंतर} \mr{पश्चात्ताप} \mr{झाला}.
\end{quote}

\begin{quote}
\emph{rāgāt bollyācā tyālā nantar paścāttāp jhālā.}
\end{quote}

\begin{quote}
Later he felt remorse for speaking in anger.
\end{quote}

Masculine.

\mr{पश्चात्}, after, + \mr{ताप}, burning: the pain that comes afterwards. \mr{पश्चात्ताप} \mr{होणे}, with the dative.

\begin{cousinweb}
\textbf{\mr{पश्चात्ताप}} (\emph{paścāttāp}) is a tatsama, taken from Sanskrit \emph{paścāttāpa}, \textquotedblleft{}after-pain, repentance\textquotedblright{} (Monier-Williams), built on \emph{tāpa}, \textquotedblleft{}heat, pain\textquotedblright{}.
\end{cousinweb}

\begin{grammarlens}[title={What you've built}]
One more word for describing what you lived and felt.
\end{grammarlens}

\subsection*{Your turn}
\begin{itemize}
\raggedright
  \item \emph{Say it:} \emph{paścāttāp}
  \item \emph{Say it:} \emph{paścāttāp}, once more
  \item \emph{Say it:} \emph{rāgāt bollyācā tyālā nantar paścāttāp jhālā.} — Later he felt remorse for speaking in anger.
  \item \emph{Recall:} say the Marathi for to be bashful, then the Marathi for regret, then say \emph{paścāttāp} again
\end{itemize}

\begin{tcolorbox}[breakable,colback=teal!4,colframe=teal!35!black,title={Before you move on}]
What does \mr{पश्चात्ताप} mean? (\textquotedblleft{}Remorse\textquotedblright{}.) Say it once more.
\end{tcolorbox}
\section[\texorpdfstring{\mr{हुंदका} (hundkā)}{हुंदका (hundkā)}]{\mr{हुंदका} (hundkā) — a sob}

\label{lesson:MR-C457-hundka}



\begin{quote}
Before the new one: say the Marathi for regret, then the Marathi for remorse.
\end{quote}

\subsection*{You'll want to know: \mr{हुंदका}}
\textbf{\mr{हुंदका}} — \emph{hundkā} — \textquotedblleft{}a sob\textquotedblright{}.

\begin{quote}
\mr{निरोप} \mr{घेताना} \mr{तिला} \mr{हुंदका} \mr{आला}.
\end{quote}

\begin{quote}
\emph{nirop ghetānā tilā hundkā ālā.}
\end{quote}

\begin{quote}
A sob caught in her throat as she said goodbye.
\end{quote}

Masculine.

\mr{हुंदका} \mr{येणे}, a sob comes, with the dative; \mr{हुंदके} \mr{देणे}, to sob.

\begin{cousinweb}
\textbf{\mr{हुंदका}} (\emph{hundkā}) is a Marathi noun whose origin is not settled.
\end{cousinweb}

\begin{grammarlens}[title={What you've built}]
One more word for describing what you lived and felt.
\end{grammarlens}

\subsection*{Your turn}
\begin{itemize}
\raggedright
  \item \emph{Say it:} \emph{hundkā}
  \item \emph{Say it:} \emph{hundkā}, once more
  \item \emph{Say it:} \emph{nirop ghetānā tilā hundkā ālā.} — A sob caught in her throat as she said goodbye.
  \item \emph{Recall:} say the Marathi for regret, then the Marathi for remorse, then say \emph{hundkā} again
\end{itemize}

\begin{tcolorbox}[breakable,colback=teal!4,colframe=teal!35!black,title={Before you move on}]
What does \mr{हुंदका} mean? (\textquotedblleft{}A sob\textquotedblright{}.) Say it once more.
\end{tcolorbox}
\section[\texorpdfstring{\mr{हळहळणे} (haḷhaḷṇe)}{हळहळणे (haḷhaḷṇe)}]{\mr{हळहळणे} (haḷhaḷṇe) — to feel sorry; to lament}

\label{lesson:MR-C457-halhalne}



\begin{quote}
Before the new one: say the Marathi for remorse, then the Marathi for a sob.
\end{quote}

\subsection*{You'll want to know: \mr{हळहळणे}}
\textbf{\mr{हळहळणे}} — \emph{haḷhaḷṇe} — \textquotedblleft{}to feel sorry; to lament\textquotedblright{}.

\begin{quote}
\mr{जुनं} \mr{झाड} \mr{तोडलं} \mr{म्हणून} \mr{आजी} \mr{हळहळली}.
\end{quote}

\begin{quote}
\emph{junaṁ jhāḍ toḍlaṁ mhaṇūn ājī haḷhaḷlī.}
\end{quote}

\begin{quote}
Grandmother felt sorry that the old tree was cut down.
\end{quote}

Past (he …): \mr{तो} \mr{हळहळला}, \emph{to haḷhaḷlā}.

Sorrow for a loss: \mr{हळहळ} \mr{वाटणे}, to feel a pang of regret.

\begin{cousinweb}
\textbf{\mr{हळहळणे}} (\emph{haḷhaḷṇe}) is a Marathi verb built on a doubled syllable; its origin is not settled.
\end{cousinweb}

\begin{grammarlens}[title={What you've built}]
One more word for describing what you lived and felt.
\end{grammarlens}

\subsection*{Your turn}
\begin{itemize}
\raggedright
  \item \emph{Say it:} \emph{haḷhaḷṇe}
  \item \emph{Say it:} \emph{haḷhaḷṇe}, once more
  \item \emph{Say it:} \emph{junaṁ jhāḍ toḍlaṁ mhaṇūn ājī haḷhaḷlī.} — Grandmother felt sorry that the old tree was cut down.
  \item \emph{Recall:} say the Marathi for remorse, then the Marathi for a sob, then say \emph{haḷhaḷṇe} again
\end{itemize}

\begin{tcolorbox}[breakable,colback=teal!4,colframe=teal!35!black,title={Before you move on}]
What does \mr{हळहळणे} mean? (\textquotedblleft{}To feel sorry; to lament\textquotedblright{}.) Say it once more.
\end{tcolorbox}
\section[\texorpdfstring{\mr{मुसमुसणे} (musmusṇe)}{मुसमुसणे (musmusṇe)}]{\mr{मुसमुसणे} (musmusṇe) — to sob quietly}

\label{lesson:MR-C457-musmusne}



\begin{quote}
Before the new one: say the Marathi for a sob, then the Marathi for to feel sorry.
\end{quote}

\subsection*{You'll want to know: \mr{मुसमुसणे}}
\textbf{\mr{मुसमुसणे}} — \emph{musmusṇe} — \textquotedblleft{}to sob quietly\textquotedblright{}.

\begin{quote}
\mr{खोलीत} \mr{बसून} \mr{मुलगी} \mr{मुसमुसत} \mr{होती}.
\end{quote}

\begin{quote}
\emph{kholīt basūn mulgī musmusat hotī.}
\end{quote}

\begin{quote}
The girl sat in the room, sobbing quietly.
\end{quote}

Past (he …): \mr{तो} \mr{मुसमुसला}, \emph{to musmuslā}.

Quieter than \mr{रडणे}: sniffs heard through a door.

\begin{cousinweb}
\textbf{\mr{मुसमुसणे}} (\emph{musmusṇe}) is inherited: CDIAL lists Marathi \emph{musmusṇẽ}, \textquotedblleft{}to sob\textquotedblright{}, under the expressive root \emph{*muss} (CDIAL 10227).
\end{cousinweb}

\begin{grammarlens}[title={What you've built}]
One more word for describing what you lived and felt.
\end{grammarlens}

\subsection*{Your turn}
\begin{itemize}
\raggedright
  \item \emph{Say it:} \emph{musmusṇe}
  \item \emph{Say it:} \emph{musmusṇe}, once more
  \item \emph{Say it:} \emph{kholīt basūn mulgī musmusat hotī.} — The girl sat in the room, sobbing quietly.
  \item \emph{Recall:} say the Marathi for a sob, then the Marathi for to feel sorry, then say \emph{musmusṇe} again
\end{itemize}

\begin{tcolorbox}[breakable,colback=teal!4,colframe=teal!35!black,title={Before you move on}]
What does \mr{मुसमुसणे} mean? (\textquotedblleft{}To sob quietly\textquotedblright{}.) Say it once more.
\end{tcolorbox}
\section[(dialogue)]{Looking back: feelings from mind and heart}

\label{lesson:MR-R457-review-450-457}



\begin{quote}
Say the Marathi for the mind and life.
\end{quote}

\subsection*{Your turn}
Say these aloud:

\begin{itemize}
\raggedright
  \item the mind, life, happiness, surprise, to be deeply moved
  \item cheerful, gloomy, pleased, displeased, to get fed up
  \item a shock, a shiver, to take fright, to panic, to be stunned
  \item hot-tempered, to be enraged, to get irritated, to snap at, to grumble
  \item a reconciliation, a misunderstanding, to sulk, to coax round, to be crestfallen
  \item dread, restless, a wearing rush, to toss and turn, to feel stifled
  \item self-confidence, hesitation, to shrink back, to get flustered, to be bashful
  \item regret, remorse, a sob, to feel sorry, to sob quietly
\end{itemize}

\begin{tcolorbox}[breakable,colback=teal!4,colframe=teal!35!black,title={Before you move on}]
The mind? (\textbf{\mr{मन}}, \emph{man}.) Life? (\textbf{\mr{जीव}}, \emph{jīv}.) Happiness? (\textbf{\mr{सुख}}, \emph{sukh}.) Surprise? (\textbf{\mr{आश्चर्य}}, \emph{āścarya}.) To be deeply moved? (\textbf{\mr{हेलावणे}}, \emph{helāvṇe}.) Cheerful? (\textbf{\mr{प्रसन्न}}, \emph{prasanna}.) Gloomy? (\textbf{\mr{उदास}}, \emph{udās}.) Pleased? (\textbf{\mr{खुश}}, \emph{khuś}.) Displeased? (\textbf{\mr{नाराज}}, \emph{nārāj}.) To get fed up? (\textbf{\mr{वैतागणे}}, \emph{vaitāgṇe}.) A shock? (\textbf{\mr{धक्का}}, \emph{dhakkā}.) A shiver? (\textbf{\mr{शहारा}}, \emph{śahārā}.) To take fright? (\textbf{\mr{धसकणे}}, \emph{dhasakṇe}.) To panic? (\textbf{\mr{बिथरणे}}, \emph{bitharṇe}.) To be stunned? (\textbf{\mr{हबकणे}}, \emph{habakṇe}.) Hot-tempered? (\textbf{\mr{तापट}}, \emph{tāpaṭ}.) To be enraged? (\textbf{\mr{संतापणे}}, \emph{santāpṇe}.) To get irritated? (\textbf{\mr{चिडणे}}, \emph{ciḍṇe}.) To snap at? (\textbf{\mr{खेकसणे}}, \emph{khekasṇe}.) To grumble? (\textbf{\mr{कुरकुरणे}}, \emph{kurkurṇe}.) A reconciliation? (\textbf{\mr{समेट}}, \emph{sameṭ}.) A misunderstanding? (\textbf{\mr{गैरसमज}}, \emph{gairsamaj}.) To sulk? (\textbf{\mr{रुसणे}}, \emph{rusṇe}.) To coax round? (\textbf{\mr{मनवणे}}, \emph{manavṇe}.) To be crestfallen? (\textbf{\mr{हिरमुसणे}}, \emph{hirmusṇe}.) Dread? (\textbf{\mr{धास्ती}}, \emph{dhāstī}.) Restless? (\textbf{\mr{बेचैन}}, \emph{becain}.) A wearing rush? (\textbf{\mr{दगदग}}, \emph{dagdag}.) To toss and turn? (\textbf{\mr{तळमळणे}}, \emph{taḷmaḷṇe}.) To feel stifled? (\textbf{\mr{घुसमटणे}}, \emph{ghusmaṭṇe}.) Self-confidence? (\textbf{\mr{आत्मविश्वास}}, \emph{ātmaviśvās}.) Hesitation? (\textbf{\mr{संकोच}}, \emph{saṅkoc}.) To shrink back? (\textbf{\mr{बिचकणे}}, \emph{bicakṇe}.) To get flustered? (\textbf{\mr{गांगरणे}}, \emph{gāṅgarṇe}.) To be bashful? (\textbf{\mr{बुजणे}}, \emph{bujṇe}.) Regret? (\textbf{\mr{खेद}}, \emph{khed}.) Remorse? (\textbf{\mr{पश्चात्ताप}}, \emph{paścāttāp}.) A sob? (\textbf{\mr{हुंदका}}, \emph{hundkā}.) To feel sorry? (\textbf{\mr{हळहळणे}}, \emph{haḷhaḷṇe}.) To sob quietly? (\textbf{\mr{मुसमुसणे}}, \emph{musmusṇe}.)
\end{tcolorbox}
